%%%%%%%%%%%%%%%%%
% This is an sample CV template created using altacv.cls
% (v1.7, 9 August 2023) written by LianTze Lim (liantze@gmail.com). Compiles with pdfLaTeX, XeLaTeX and LuaLaTeX.
%
%% It may be distributed and/or modified under the
%% conditions of the LaTeX Project Public License, either version 1.3
%% of this license or (at your option) any later version.
%% The latest version of this license is in
%%    http://www.latex-project.org/lppl.txt
%% and version 1.3 or later is part of all distributions of LaTeX
%% version 2003/12/01 or later.
%%%%%%%%%%%%%%%%

%% Use the "normalphoto" option if you want a normal photo instead of cropped to a circle
% \documentclass[10pt,a4paper,normalphoto]{altacv}

\documentclass[10pt,a4paper,ragged2e,withhyper]{altacv}
%% AltaCV uses the fontawesome5 and packages.
%% See http://texdoc.net/pkg/fontawesome5 for full list of symbols.

% Change the page layout if you need to
\geometry{left=1.25cm,right=1.25cm,top=1.5cm,bottom=1.5cm,columnsep=0cm}

% The paracol package lets you typeset columns of text in parallel
\usepackage{paracol}

% Change the font if you want to, depending on whether
% you're using pdflatex or xelatex/lualatex
% WHEN COMPILING WITH XELATEX PLEASE USE
% xelatex -shell-escape -output-driver="xdvipdfmx -z 0" sample.tex
\ifxetexorluatex
  % If using xelatex or lualatex:
  \setmainfont{Roboto Slab}
  \setsansfont{Lato}
  \renewcommand{\familydefault}{\sfdefault}
\else
  % If using pdflatex:
  \usepackage[rm]{roboto}
  \usepackage[defaultsans]{lato}
  % \usepackage{sourcesanspro}
  \renewcommand{\familydefault}{\sfdefault}
\fi

% Change the colours if you want to
\definecolor{SlateGrey}{HTML}{2E2E2E}
\definecolor{LightGrey}{HTML}{666666}
\definecolor{DarkPastelRed}{HTML}{450808}
\definecolor{PastelRed}{HTML}{8F0D0D}
\definecolor{GoldenEarth}{HTML}{E7D192}
\colorlet{name}{black}
\colorlet{tagline}{PastelRed}
\colorlet{heading}{DarkPastelRed}
\colorlet{headingrule}{GoldenEarth}
\colorlet{subheading}{PastelRed}
\colorlet{accent}{PastelRed}
\colorlet{emphasis}{SlateGrey}
\colorlet{body}{LightGrey}

% Change some fonts, if necessary
\renewcommand{\namefont}{\Huge\rmfamily\bfseries}
\renewcommand{\personalinfofont}{\footnotesize}
\renewcommand{\cvsectionfont}{\LARGE\rmfamily\bfseries}
\renewcommand{\cvsubsectionfont}{\large\bfseries}


% Change the bullets for itemize and rating marker
% for \cvskill if you want to
\renewcommand{\cvItemMarker}{{\small\textbullet}}
\renewcommand{\cvRatingMarker}{\faCircle}
% ...and the markers for the date/location for \cvevent
% \renewcommand{\cvDateMarker}{\faCalendar*[regular]}
% \renewcommand{\cvLocationMarker}{\faMapMarker*}


% If your CV/résumé is in a language other than English,
% then you probably want to change these so that when you
% copy-paste from the PDF or run pdftotext, the location
% and date marker icons for \cvevent will paste as correct
% translations. For example Spanish:
% \renewcommand{\locationname}{Ubicación}
% \renewcommand{\datename}{Fecha}


%% Use (and optionally edit if necessary) this .tex if you
%% want to use an author-year reference style like APA(6)
%% for your publication list
% \input{pubs-authoryear.tex}

%% Use (and optionally edit if necessary) this .tex if you
%% want an originally numerical reference style like IEEE
%% for your publication list
\input{pubs-num.tex}

%% sample.bib contains your publications
\addbibresource{sample.bib}

\begin{document}
\name{Myoungchul Kim}
\tagline{Data Scientist}
%% You can add multiple photos on the left or right
%\photoR{2.8cm}{Globe_High}
% \photoL{2.5cm}{Yacht_High,Suitcase_High}

\personalinfo{%
  % Not all of these are required!
  \lang{English: Fluent}
  \lang{Japanese: Fluent, JLPT N1}
  \lang{Korean: Native}\\
  \smallskip
  \email{kmc8907@gmail.com}
  \phone{080-9801-7956}
  \location{Greater Tokyo Area, Japan}
  \github{MyoungchulK}
  \linkedin{myoungchulkim}
  \homepage{troopl.com/kmc8907}
  \orcid{0000-0002-8624-5564}
  %% You can add your own arbitrary detail with
  %% \printinfo{symbol}{detail}[optional hyperlink prefix]
  % \printinfo{\faPaw}{Hey ho!}[https://example.com/]

  %% Or you can declare your own field with
  %% \NewInfoFiled{fieldname}{symbol}[optional hyperlink prefix] and use it:
  % \NewInfoField{gitlab}{\faGitlab}[https://gitlab.com/]
  % \gitlab{your_id}
  %%
  %% For services and platforms like Mastodon where there isn't a
  %% straightforward relation between the user ID/nickname and the hyperlink,
  %% you can use \printinfo directly e.g.
  % \printinfo{\faMastodon}{@username@instace}[https://instance.url/@username]
  %% But if you absolutely want to create new dedicated info fields for
  %% such platforms, then use \NewInfoField* with a star:
  % \NewInfoField*{mastodon}{\faMastodon}
  %% then you can use \mastodon, with TWO arguments where the 2nd argument is
  %% the full hyperlink.
  % \mastodon{@username@instance}{https://instance.url/@username}
}

\makecvheader
%% Depending on your tastes, you may want to make fonts of itemize environments slightly smaller
% \AtBeginEnvironment{itemize}{\small}

%% Set the left/right column width ratio to 6:4.
\columnratio{1}

% Start a 2-column paracol. Both the left and right columns will automatically
% break across pages if things get too long.
%\begin{paracol}{2}
\cvsection{About Me}
\begin{quote}
Data Scientist with a strong foundation in {\bf AstroParticle Physics}, {\bf statistical modeling}, and {\bf large-scale data analysis}. Experienced in {\bf LLM, {\bf RAG}, {\bf LangGraph}, multi-agent systems}, {\bf SFT}, {\bf RLHF}, {\bf ML pipeline} automation, and {\bf Python}-based {\bf ML Engineering}. Proven track record across academic research and industry AI applications.
\end{quote}

\medskip

\cvsection{Technical Skills}

\cveventv{Coding Tools: }
\cvtag{Python}
\cvtag{C++}
\cvtag{Git / GitHub}
\cvtag{Vim}
\cvtag{Jupyter}
\cvtag{Latex}%\\
%\quad\quad\quad\quad\quad\qquad
\cvtag{ROOT}
\cvtag{HTCondor}
\cvtag{CVMFS}

\dividerv\smallskip

\cveventv{Data Analytics: }
\cvtag{NumPy}
\cvtag{SciPy}
\cvtag{Scikit-learn}
\cvtag{Pandas}
\cvtag{GeoPandas}
\cvtag{SQL}
\cvtag{Matplotlib}
\cvtag{Seaborn}

\dividerv\smallskip

\cveventv{Modelling: }
\cvtag{TensorFlow}
\cvtag{PyTorch}
\cvtag{Unsupervised Leaning}
\cvtag{CNN}
\cvtag{GNN}
%\quad\quad\quad\quad\quad\qquad
\cvtag{Time series}
\cvtag{Ensemble Methods}\\
\quad\quad\quad\quad\quad\qquad
\cvtag{Statsmodels}
\cvtag{LLM Agent}
\cvtag{RAG}
\cvtag{RLHF}
\cvtag{SFT}
\cvtag{Hugging Face}

\dividerv\smallskip

\cveventv{Deployment: }
\cvtag{AWS}
\cvtag{GCP}
\cvtag{Docker}
\cvtag{Databricks}
\cvtag{Dagster}
\cvtag{LangChain}
\cvtag{LangGraph}
\cvtag{FastAPI}
\cvtag{Streamlit}\\
\quad\quad\quad\quad\quad\qquad
\cvtag{Raspberry Pi}

\medskip
%\smallskip

\cvsection{Experience}

\cvevent{Data Scientist}{Fracta}{November 2024 -- Ongoing}{Tokyo, Japan}{}{}

\cvsubevent{AI-Driven Infrastructure Assessments}{November 2024 -- Ongoing}{Role: Data Scientist}{Team size: 4}
\begin{itemize}
\item Developed an {\bf LLM}-powered {\bf multi-agent system} leveraging {\bf LangGraph} and {\bf RAG} for automated raw data preprocessing prior to {\bf ML} ingestion.
\item Deployed {\bf ML models} on large-scale infrastructure networks to predict {\bf Likelihood of Failure (LOF)} and deliver client-facing risk assessments.
\item Engineered end-to-end {\bf ML pipeline} automation, improving scalability and reducing manual intervention.
\end{itemize}
\divider

\cvevent{Data Science Instructor}{Le Wagon}{October 2024 -- Ongoing}{Tokyo, Japan}{}{}
\cvsubevent{Data Science \& AI Bootcamp Instruction}{October 2024 -- Ongoing}{Role: Teaching Assistant}{Team size: 13}
\begin{itemize}
\item Instructed {\bf Python} for Data Science, covering data extraction, manipulation, visualization, {\bf and statistics}.
\item Guided students through {\bf ML/DL} workflows using {\bf Scikit-Learn} and {\bf neural network} architecture design.
\item Delivered training in {\bf ML Engineering}, including {\bf Python} package deployment on {\bf GCP} and ethical implications of {\bf AI}.
\end{itemize}
\divider

\cvevent{Python Data Scientist/Analyst}{Turing}{July 2024 -- January 2025}{Remote}{}{}

\cvsubevent{Multilingual SFT and RLHF Implementation for Next-Generation AI}{July 2024 -- January 2025}{Role: LLM Dataset Creation}{Team size: 6}
\begin{itemize}
\item Performed {\bf Supervised Fine-Tuning (SFT)} and {\bf Reinforcement Learning with Human Feedback (RLHF)} across {\bf English, Korean}, and {\bf Japanese} datasets to improve {\bf LLM} alignment.
\item Analyzed model performance via {\bf Loss Pattern} tracking and {\bf Output-Based Evaluation} to enhance robustness and reliability.
\item Curated model responses through {\bf Comparative Analysis} and {\bf Preference Pair} evaluation, ensuring contextually accurate outputs.
\end{itemize}
\divider

\cvevent{AI Training for Japanese Writers}{Outlier}{July 2024 -- October 2024}{Remote}{}{}

\cvsubevent{Fine-Tuning and Validation of AI-Generated Japanese Voice}{July 2024 -- October 2024}{Role: LLM evaluation}{Team size: 12}
\begin{itemize}
\item Validated {\bf AI-generated voice} outputs, assessing quality via {\bf data augmentation, pitch}, and {\bf speech rate} for natural and contextually accurate results.
\item Applied {\bf linguistic expertise} in Japanese to evaluate {\bf voice model} performance against clarity, consistency, and {\bf cultural accuracy} benchmarks.
\end{itemize}
\divider

\cvevent{Graduate Research Assistant, Ph.D.}{International Center for Hadron Astrophysics (ICEHAP)}{April 2017 -- December 2023}{Chiba University, Japan}{}{}

\cvsubevent{\href{https://user-web.icecube.wisc.edu/~mkim/ARA_MF/index.html}{Search for Ultra-high Energy Neutrinos from Askaryan Radio Array (ARA) by Template Method \faLink}}{April 2018 -- December 2023}{Role: Project Leader}{Team size: 6}
\begin{itemize}
\item Classified astronomical signals using {\bf PCA} on ~200 TB of {\bf radio-frequency} data, optimized via {\bf Frequentist Statistics} and {\bf Pseudo Experiments}.
\item Implemented {\bf FFT}, {\bf Interferometry}, and {\bf Matched Filter} for feature extraction on large-scale {\bf CPU \& GPU} clusters using scientific {\bf Python} and {\bf C++}.
\item Evaluated results via {\bf statistical significance}, {\bf systematic uncertainty}, and {\bf Monte Carlo simulation}; pipeline adopted by international collaborators.
\end{itemize}
\cvsubevent{Development of In Situ Antenna Model for Simulation}{April 2017 -- March 2019}{Role: Data Management}{Team size: 4}
\begin{itemize}
\item Performed high-precision calibration for {\bf radio-frequency antennas} to measure and characterize physical properties.
\item Extracted feature patterns from raw data through {\bf noise \& signal analysis}, contributing to simulation model development.
\end{itemize}
\divider

\cvevent{Graduate Research Assistant, MSc.}{Neutrino AstroParticle Physics Lab (NAPPL)}{March 2015 -- February 2017}{SungKyunKwan University, South Korea}{}{}
\cvsubevent{\href{https://pos.sissa.it/236/1145/pdf}{IceCube Camera System to Study Properties of the Antarctic Ice \faLink}}{March 2015 -- February 2017}{Role: Project Leader}{Team size: 5}
\begin{itemize}
\item Classified intrinsic camera noise under extreme low-temperature conditions using {\bf feature extraction} and {\bf image analysis}.
\item Developed {\bf Python} package for camera control via {\bf Raspberry Pi}, enabling automated data collection.
\item Evaluated camera system performance by applying {\bf statistical techniques} to large-scale {\bf image datasets}.
\end{itemize}
%\divider

%\cvevent{Japanese to English \& Korean Translator}{Wovn.io}{February 2018 -- August 2018}{Tokyo}{Role: Translator}{Team size: 5}
%\begin{itemize}
%\item Provided translation services and real-time deployment for the client company website written in Japanese, incorporating cultural nuances to enhance readability and accessibility.
%\end{itemize}
\medskip

\cvsection{Project}

\cvevent{\href{https://github.com/MyoungchulK/Sound_To_Symphony_LeWgon}{AI Music Generation via RNN \faLink}}{Le Wagon Data Science \& AI Bootcamp}{January 2024 -- March 2024}{Le Wagon, Tokyo}{Role: Data Management}{Team size: 4}
\begin{itemize}
\item Architected {\bf RNN} to learn musical patterns from large classical music datasets encoded in numerical form, enabling customizable music generation.
\item Deployed end-to-end pipeline via {\bf FastAPI} and {\bf Streamlit}, integrating generated outputs with {\bf Ableton}.
\item Managed data preprocessing and model training workflows, ensuring scalable and reproducible {\bf deep learning} pipelines.
\end{itemize}
%\medskip
% use ONLY \newpage if you want to force a page break for
% ONLY the current column
%\newpage

\cvsection{Education}

\cvevent{Data Science \& AI Bootcamp}{Le Wagon}{January 2024 -- March 2024}{Japan}{}{}
\begin{itemize}
\item Studied {\bf Python} for Data Science, covering {\bf data extraction}, {\bf manipulation}, {\bf visualization}, {\bf statistics}, and {\bf linear algebra}.
\item Built end-to-end {\bf ML/DL} workflows using {\bf Scikit-Learn} and designed {\bf neural network} architectures for practical applications.
\item Developed {\bf Python} packages for large-scale data tasks on {\bf GCP}, with exposure to {\bf ML Engineering} and {\bf AI ethics}.
\end{itemize}
\divider

\cvevent{Completion of Ph.D. program (ABD), AstroParticle Physics}{Chiba University}{April 2017 -- December 2023}{Japan}{}{}
\begin{itemize}
\item Research topic: Search for Ultra-high Energy Neutrinos Using Eight Years of Data from Two ARA Stations by the Neutrino Template Method
\end{itemize}
\divider

\cvevent{Master of Science, AstroParticle Physics}{SungKyunKwan University}{March 2015 -- February 2017}{South Korea}{}{}
\begin{itemize}
\item Thesis title: Performance study of camera system for the IceCube-Gen2 detector
\end{itemize}
\divider

\cvevent{Bachelor in Science, Physics}{SungKyunKwan University}{March 2011 -- February 2015}{South Korea}{}{}

%\switchcolumn




%\cvsection{Languages}
%\cvskillv{English}{Fluent}\\
%\divider
%\cvskillv{Japanese}{JLPT N1, Fluent}\\ %% Supports X.5 values.
%\divider
%\cvskillv{Korean}{Native}\\

%% Yeah I didn't spend too much time making all the
%% spacing consistent... sorry. Use \smallskip, \medskip,
%% \bigskip, \vspace etc to make adjustments.



% Adapted from @Jake's answer from http://tex.stackexchange.com/a/82729/226
% \wheelchart{outer radius}{inner radius}{
% comma-separated list of value/text width/color/detail}
%\wheelchart{1.5cm}{0.5cm}{%
%  %8/8em/accent!60/Contrabass,
%  2/10em/accent/Orchestra,
%  5/6em/accent!20/Classical Music
%}

% \divider

%\newpage

\cvsection{Selected Publications}

%% Specify your last name(s) and first name(s) as given in the .bib to automatically bold your own name in the publications list.
%% One caveat: You need to write \bibnamedelima where there's a space in your name for this to work properly; or write \bibnamedelimi if you use initials in the .bib
%% You can specify multiple names, especially if you have changed your name or if you need to highlight multiple authors.
\mynames{Lim/Lian\bibnamedelima Tze,
  Wong/Lian\bibnamedelima Tze,
  Lim/Tracy,
  Lim/L.\bibnamedelimi T.}
%% MAKE SURE THERE IS NO SPACE AFTER THE FINAL NAME IN YOUR \mynames LIST

\nocite{*}

\printbibliography[heading=pubtype,title={\printinfo{\faFile*[regular]}{Journal Articles}},type=article]


\cvsection{Awards}

\cvachievement{\faTrophy}{Japanese Government Monbukagakusho Scholarship (MEXT)}{Graduate Research Assistant in Ph.D., 2017 -- 2020}
\divider

\cvachievement{\faHeartbeat}{Teaching Assistant (T.A.) of Korea \& Japan Joint Government Scholarship}{Teaching Assistant for a freshman Korean students, 2017 -- 2018}
\divider

\cvachievement{\faTrophy}{BK21+ Research student scholarship}{Graduate Research Assistant in Msc., 2015 -- 2017}

\divider

\cvachievement{\faHeartbeat}{Operating Assistant scholarship}{Physics Experiment Assistant, 2016 -- 2017}

\divider

\cvachievement{\faTrophy}{CK Research student scholarship}{Research Assistant in Bsc., 2014}


\cvsection{Interests}

\cvtag{Classical Music}
\cvtag{Orchestra}
\cvtag{Contrabass}
\cvtag{Universe}
\cvtag{Fourier Transform}

%\end{paracol}


\end{document}
